%==============================================================================
% sections/01_introduction.tex
%==============================================================================
\section{Introduction}
\label{sec:intro}

\subsection{Motivation}

Heavy-tailed risk models are built around an asymmetry between ordinary
fluctuations and rare losses. When the loss distribution is regularly varying,
large portfolio exceedances are not well represented by Gaussian perturbations:
they are typically generated by one extreme component, one extreme latent shock,
one extreme block, or one extreme vector direction. The simplest version of this
idea is the single-big-jump principle for independent heavy-tailed summands.
Financial factor portfolios, however, rarely contain exactly independent factors.
Market, size, value, profitability, investment, country, sector, credit,
volatility, liquidity, and residual terms can co-move in their extremes. This
paper therefore develops the independent theory carefully and then turns it into
a staged framework for non-independent factor-loss models.

Let
\[
  S_N=\sum_{i=1}^N X_i
\]
be the tail-risk target, where the $X_i$ are signed one-sided loss contributions
obtained from factor returns, portfolio exposures, and residuals. The independent
baseline assumes
\[
  \Pb(X_i>x)\sim c_i\Gb(x),\qquad \Gb\in \RV_{-\alpha},
\]
for a common reference tail and a finite number of summands. The leading question
is the rare-event probability
\[
  \mu(x)=\Pb(S_N>x),\qquad x\to\infty.
\]
The computational question is how to estimate $\mu(x)$ accurately when crude
Monte Carlo sees few or no exceedances.

The paper has two complementary objectives. First, it sharpens the independent
finite-$N$ theory and the conditional Monte Carlo estimators used by
\citet{PourbabaeeSolari2019}. Second, it develops a concrete route from that
independent theory to non-independent factor models. The dependent route is not
presented as a single over-general theorem. Instead, it is organized by the
source of tail dependence: arbitrary conditional tails, latent independent
shocks, common shocks, independent blocks, copulas, multivariate regular
variation, and hidden regular variation.

\subsection{Core message}

The core message is a hierarchy of rare-event mechanisms:
\[
  \text{one large coordinate}
  \subset
  \text{one large latent shock}
  \subset
  \text{one large block}
  \subset
  \text{one large spectral vector shock}
  \subset
  \text{one large hidden-cone shock}.
\]
The independent theory is the coordinate-axis special case. A common-shock
factor model replaces several observed extreme coordinates by a single latent
source. A block model permits multiple components to move together inside a
sector, region, or factor cluster. A multivariate regularly varying model
replaces the active coordinate set by a spectral measure over extreme directions.
Hidden regular variation adds slower joint-tail cones that are invisible to the
first-order limit but can dominate ordinary second-order marginal corrections.

\begin{figure}[t]
\centering
\begin{tikzpicture}[
  node distance=0.8cm and 0.9cm,
  box/.style={rectangle, rounded corners, draw, align=center, minimum width=3.2cm, minimum height=0.9cm, font=\small},
  arr/.style={->, thick, >=Stealth}
]
\node[box, fill=blue!10] (coord) {One large\\observed coordinate};
\node[box, fill=green!12, right=of coord] (latent) {One large\\latent shock};
\node[box, fill=yellow!18, right=of latent] (radial) {One large radial\\vector shock};
\node[box, fill=orange!15, right=of radial] (hidden) {One large hidden\\pair/block cone};
\draw[arr] (coord) -- (latent);
\draw[arr] (latent) -- (radial);
\draw[arr] (radial) -- (hidden);
\node[align=center, font=\scriptsize, below=0.35cm of coord] {independent CdMC};
\node[align=center, font=\scriptsize, below=0.35cm of latent] {shock-basis CdMC};
\node[align=center, font=\scriptsize, below=0.35cm of radial] {spectral CdMC};
\node[align=center, font=\scriptsize, below=0.35cm of hidden] {mixture cone CdMC};
\end{tikzpicture}

\caption[Hierarchy of independent and dependent big-jump mechanisms]{
  Hierarchy of big-jump mechanisms developed in the paper. The independent
  finite-$N$ theory corresponds to axis mass, latent-shock models move the
  conditioning to independent underlying shocks, block models allow within-block
  joint extremes, multivariate regular variation integrates over a spectral
  measure, and hidden regular variation captures slower joint-tail cones that may
  be material in finite samples.
}
\label{fig:dependent-hierarchy}
\end{figure}

\subsection{Roadmap of assumptions}
\label{subsec:assumption-roadmap}

The theoretical results in this paper layer assumptions cumulatively rather
than invoking a single global hypothesis. The independent baseline of
\cref{sec:independent-baseline} requires only common-reference regular
variation, a non-empty active set, a deterministic tie-breaking rule, and
measurability \cref{ass:rv-common-ref,ass:active-nonempty,ass:tie,ass:measurability}.
The second-order expansion adds direct second-order tail equivalence
\cref{ass:sote} and a finite first moment ($\alpha>1$). The dependent
conditional Monte Carlo identity of \cref{sec:dependent-cdmcs} relaxes
independence to arbitrary dependence and replaces the marginal survival by the
conditional kernel; bounded-relative-error guarantees in that section add
model-specific conditions stated where they bind. The MRV extension of
\cref{sec:mrv-spectral} adds multivariate regular variation of the loss-vector
plus boundary continuity of the linear loss set under the limit measure. The
hidden-regular-variation extension of \cref{sec:hidden-rv} adds a hidden-cone
second-order assumption. \Cref{appx:notation} consolidates the symbol
inventory.

\subsection{Contributions}

\paragraph{Independent finite-$N$ asymptotics.}
\Cref{sec:independent-baseline} proves the exact finite-portfolio catastrophe
principle for the maximum and the corresponding single-big-jump equivalence for
the sum. For the active set $\Aset=\{i:c_i>0\}$, the leading term is
\[
  \Pb(S_N>x)\sim \sum_{i\in\Aset}\Pb(X_i>x)
  \sim \Big(\sum_{i\in\Aset}c_i\Big)\Gb(x).
\]
The section also records finite-sample bounds that make the asymptotic constants
usable as implementation checks.

\paragraph{Corrected second-order expansion.}
Under direct second-order tail equivalence, local-shift regularity, and finite
first moments, \Cref{sec:independent-baseline} derives
\[
  \Pb(S_N>x)
  = C\Gb(x)
    + \Gb(x)A(x)\sum_{i\in\Aset}c_i d_i
    + {\alpha\Gb(x)\over x}\sum_{i\in\Aset}c_i\mu_{-i}
    + o\{\Gb(x)(a(x)\vee x^{-1})\},
\]
where $C=\sum_{i\in\Aset}c_i$, $a(x)=|A(x)|$, and
$\mu_{-i}=\sum_{j\ne i}\E X_j$. The leave-one-out mean is essential: once
coordinate $i$ is the large component, the remaining portfolio supplies the
local shift.

\paragraph{Dependence extension path.}
\Cref{sec:dependent-cdmcs} proves that the same maximum-index partition gives an
exact unbiased estimator under arbitrary dependence after replacing the marginal
survival $\Fb_i$ by the conditional survival $\Pb(X_i>t\mid X_{-i})$. It then
gives implementable specializations for latent independent shocks, common shocks,
independent blocks, copulas, and vine copulas. \Cref{sec:mrv-spectral} develops
the multivariate-regular-variation extension in which the one-large-coordinate
constant becomes a spectral-measure integral. \Cref{sec:hidden-rv} formulates
hidden regular variation as a second-order joint-tail correction when ordinary
MRV puts first-order mass on the coordinate axes but slower cones remain
material.

\paragraph{Estimator family and efficiency criteria.}
\Cref{sec:estimators-efficiency} compares crude Monte Carlo, independent CdMC,
dependent conditional-tail CdMC, latent-shock CdMC, block CdMC, spectral/radial
CdMC, and control-variate variants. In the independent case, the summed CdMC
estimator
\[
  Z^{\CdMC}(x)=\sum_{i=1}^N \Fb_i((x-S_{-i})\vee M_{-i})
\]
is unbiased and has the scale-invariant bounded-relative-error bound
$N^\alpha-1$. In dependent settings, the same partition gives unbiasedness, but
bounded-relative-error guarantees require conditional-tail, block, latent-shock,
or radial-envelope assumptions.

\begin{figure}[t]
\centering
% D3: Estimator flowchart
\begin{tikzpicture}[
  est/.style={rectangle, draw, rounded corners, minimum width=2.55cm, minimum height=0.9cm, align=center, font=\small},
  lab/.style={font=\scriptsize, fill=white, inner sep=1pt},
  >=Stealth, node distance=1.15cm and 1.85cm
]
\node[est, fill=red!10] (crmc) {CrMC \\ \footnotesize baseline};
\node[est, fill=yellow!20, right=of crmc] (cdmc) {Summed CdMC \\ \footnotesize unbiased, BRE};
\node[est, fill=yellow!20, above right=of cdmc] (strat) {Stratified CdMC \\ \footnotesize one index, BRE};
\node[est, fill=blue!10, right=of cdmc] (ana) {Pareto/Lomax \\ \footnotesize closed form};
\node[est, fill=green!15, below right=of cdmc] (vre) {Control variate \\ \footnotesize exact centering, VRE};
\draw[->, shorten >=2pt, shorten <=2pt] (crmc) -- node[lab, above]{condition} (cdmc);
\draw[->, shorten >=2pt, shorten <=2pt] (cdmc) -- node[lab, above, sloped]{sample index} (strat);
\draw[->, shorten >=2pt, shorten <=2pt] (cdmc) -- node[lab, above]{closed form} (ana);
\draw[->, shorten >=2pt, shorten <=2pt] (cdmc) -- node[lab, below, sloped]{centered proxy} (vre);
\end{tikzpicture}

\caption[Estimator family]{
  Conditional Monte Carlo estimator family. Summed CdMC analytically conditions
  on each candidate dominant summand. Stratified CdMC samples one conditioning
  index per replicate. Pareto/Lomax closed form is a specialization of CdMC, and
  the centered control-variate path uses exact or sample-split centering.
}
\label{fig:estimators-flow}
\end{figure}

\paragraph{Open-source reference implementation.}
All estimators, diagnostics, and reproducibility utilities developed in this
paper are implemented in the \texttt{FactorTail} Python library, available at
\url{https://github.com/osolari/FactorTail}. The library mirrors the
manuscript's module structure: independent CdMC, dependent kernel CdMC,
latent-shock CdMC, block CdMC, spectral/radial CdMC, hidden-cone diagnostics,
and the full Fama--French / CRSP rolling VaR/ES pipeline. Section references
in the manuscript correspond to module paths in the library, and every
algorithm in \cref{appx:algorithms} has a one-to-one Python entry point.

\paragraph{Long-form empirical design (planned).}
\Cref{sec:simulation-plan} and \cref{sec:real-data} are written as complete
protocols rather than short placeholders. The simulation design specifies
data-generating processes, estimators-to-test, replication counts, success
criteria, and reproducibility metadata for independent and dependent regimes.
The real-data design specifies Fama--French factor data, CRSP returns, rolling
factor estimation, one-sided tail fitting, dependence diagnostics, empirical
spectral estimation, hidden-cone diagnostics, VaR/ES forecasting, estimator
runtime and variance comparisons, and backtesting. Figures and tables in
\cref{sec:simulation-plan} and \cref{sec:real-data} are implementation
placeholders with planned values marked \xx{}; numerical performance and
backtest outcomes are not asserted until generated outputs replace those
entries.

\subsection{Scope and status of empirical claims}

The theorem sections establish mathematical claims under explicit assumptions.
The simulation and real-data sections specify experiments, diagnostics, schemas,
and replacement rules. Numeric performance, backtest success, and production
runtime improvements are not asserted until generated outputs replace the \xx{}
entries. This separation is deliberate: the paper can be read as a theory paper,
an estimator paper, and an empirical protocol without confusing planned results
with observed results. The \texttt{FactorTail} library
(\url{https://github.com/osolari/FactorTail}) supplies the implementation
against which any future numeric results must be reproducible.

\subsection{Organization}

\Cref{sec:literature} places the paper in the literature. \Cref{sec:setup}
introduces notation, assumptions, and the financial loss mapping.
\Cref{sec:independent-baseline} gives the leading and second-order independent
asymptotics. \Cref{sec:dependent-cdmcs} develops exact dependent conditional
Monte Carlo identities and practical dependence modules. \Cref{sec:mrv-spectral}
and \cref{sec:hidden-rv} develop spectral and hidden-cone extensions.
\Cref{sec:estimators-efficiency} compares estimators and efficiency guarantees.
\Cref{sec:simulation-plan} specifies the planned synthetic experiments.
\Cref{sec:real-data} gives the extensive real-data analysis plan.
\Cref{sec:discussion} concludes. Appendices collect notation, proofs,
pseudo-code, reproducibility rules, dataset specifications, and experiment
manifests.
